% =====================================================================
% Parte inicial do fonte RISTI / ICITS'27 (versão em português, revisão
% duplo-cega: sem autoria): título, resumo e palavras-chave em português e
% em inglês, dentro dos limites da RISTI (150 palavras, 5 palavras-chave).
% O corpo do artigo vem de paper_pt.tex. Use sync_risti.py para gerar
% paper_pt_risti.tex e paper_pt_risti_15p.tex, e build_risti_docx.py para
% gerar o .docx. Não é compilável com LaTeX.
% =====================================================================

\title{Uma Análise por Fluxo de Acesso em Redes Viárias Dirigidas: um Estudo em Grafos de Três Bairros do Rio de Janeiro}

\abstractpt{Modelamos as redes viárias de Urca, Copacabana e Botafogo como grafos
dirigidos do OpenStreetMap. Perguntamos se indicadores estruturais ou um
fluxo de acesso identificam um bairro com uma única rota de saída, e se
isso reflete a geografia ou o tamanho. Pontes e articulações fortes e o
corte mínimo entre dois nós não ordenam os bairros. O fluxo máximo do
núcleo do bairro até as portas administrativas ordena: uma rota disjunta por ruas
na Urca, dezoito em Copacabana e quinze em Botafogo, estável sob
variações de recorte e margem. Contar interseções em vez de ruas quase
não muda os valores, e respeitar o sentido das vias reduz o acesso das
malhas abertas a menos da metade, sem alterar a Urca. Pedaços das malhas
abertas do tamanho da Urca não reproduzem o padrão, associado à geografia
da península. A contagem pode ajudar a sinalizar bairros com uma só saída.}

\keywordspt{Fluxo máximo; Mobilidade urbana; Grafos dirigidos; Conectividade forte; Redes viárias}

\entitle{An Access-Flow Analysis of Directed Street Networks: A Graph Study of Three Rio de Janeiro Neighbourhoods}

\abstracten{We model the street networks of Urca, Copacabana and Botafogo as directed
graphs built from OpenStreetMap. We ask whether structural indicators or
an access flow identify a neighbourhood with a single exit route, and
whether this reflects geography or size. Strong bridges, strong
articulation points and the node-to-node minimum cut do not order the
neighbourhoods. The maximum flow from each core to its administrative
gateways does: one street-disjoint route for Urca, eighteen for
Copacabana and fifteen for Botafogo, stable under clip and margin
variations. Counting intersections instead of streets barely changes the
values, and respecting one-way directions cuts the open grids' access to
less than half without changing Urca. Pieces of the open grids the size of
Urca do not reproduce the pattern, which is tied to the peninsula's
geography. The count can help flag neighbourhoods with a single exit.}

\keywordsen{Maximum flow; Urban mobility; Directed graphs; Strong connectivity; Street networks}

% =====================================================================
